\chapter{Szablony doby}
\markboth{Szablony doby}{Szablony doby}

Szablon to doba rozpisana na bloki, których suma daje dokładnie dwadzieścia cztery godziny. Książka ma \val{week.templates} szablony. W każdym z nich sen trwa \val{week.A.sleep} godzin, bo na nim opiera się cały plan. Wolne wynosi w szablonie A \val{week.A.free} godziny, w szablonie B \val{week.B.free} godziny, a w szablonie C \val{week.C.free} godzin. Różnią się tym, ile zostaje na pracę, dojazd i dom. Sprawdzenie sumy jest bramką narzędzia: szablon, który nie daje dwudziestu czterech godzin, nie przechodzi kontroli.

\generated{week-templates}

Zwróć uwagę na wolne. W szablonie roboczym zostaje pół godziny albo ćwierć godziny, a w weekendowym pięć. To jest miejsce, w którym większość planów pęka: wolne jest najmniej chronione, bo nie wygląda jak zobowiązanie. Dlatego w każdym szablonie jest wpisane świadomie, a nie jako reszta.

\subsection*{Jak wypełnić własny szablon}

Zacznij od bloku snu, bo w tym planie nie podlega on negocjacjom. Wpisz godzinę zasypiania i godzinę budzenia, a różnicę zapisz w godzinach. Jeśli Twoja potrzeba to dwanaście godzin, tyle właśnie zostaje. Dalej postępuj tak:

\begin{enumerate}
\item Wpisz sen jako różnicę godzin, a nie jako kwadranse wieczorem. Czterdzieści pięć minut to 0,75 godziny.
\item Zapisz realne bloki: pracę razem z przerwami, dojazd w obie strony, posiłki, ruch i higienę, dom i chleb w czasie aktywnym. Nie wpisuj tego, co chciałbyś robić, tylko to, co faktycznie zajmuje czas.
\item Zsumuj wszystko. Suma musi wynosić dokładnie 24.
\item Jeśli suma przekracza 24, skróć blok najbardziej elastyczny, a nie ten, który najbardziej lubisz. Najbardziej elastyczny to ten, którego skrócenie nie łamie żadnej innej części planu.
\item Wolne wpisz jako wartość, której bronisz, a nie jako resztę. Jeśli reszta wynosi zero, plan jest pusty w tym miejscu. Zapisz to tak, jak jest.
\end{enumerate}

\begin{center}\small
\renewcommand{\arraystretch}{1.3}
\begin{tabular}{|p{4.6cm}|p{2.2cm}|p{2.6cm}|p{4.6cm}|}
\hline
\textsf{\textbf{Blok}} & \textsf{\textbf{Godziny}} & \textsf{\textbf{Od--do}} & \textsf{\textbf{Co zastępuje, uwagi}}\\
\hline
Sen & & & \\
\hline
Praca & & & \\
\hline
Dojazd & & & \\
\hline
Ruch i higiena & & & \\
\hline
Posiłki & & & \\
\hline
Dom & & & \\
\hline
Chleb (czas aktywny) & & & \\
\hline
Wolne & & & \\
\hline
Inne & & & \\
\hline
\textbf{Razem} & \textbf{24} & & \\
\hline
\end{tabular}
\end{center}

\subsection*{Pytania kontrolne}

\begin{enumerate}
\item Czy suma Twojego szablonu wynosi dokładnie 24 godziny, czy tylko w przybliżeniu? Jeśli w przybliżeniu, który blok jest niedokładny?
\item Który blok w Twoim szablonie jest najbardziej elastyczny i co dokładnie zostaje, gdy skrócisz go o godzinę?
\item Ile godzin wolnego zostaje w dniu roboczym po odjęciu pracy i dojazdu? Policz to także w ujęciu tygodniowym.
\item Które bloki są czasem biernym, a które wymagają Twojej obecności? Czy oba rodzaje mieszczą się w tej samej godzinie?
\item Co zmienisz w szablonie, jeśli sen przesunie się o godzinę w jedną stronę, a tydzień pracy zostanie taki sam?
\end{enumerate}
